\begin{frame}{A relativistic cruise}
Eddie has always been his grandmother's favorite grandchild, so he
gets invited with her on an interstellar cruise. After the ship has accelerated
to 30\%
of the speed of light, Eddie unbuckles his seatbelt, hops out of his acceleration
couch, and extracts a folding meter-stick from his suitcase. Using the meter-stick,
he measures the length of the cabin from front to back.

\vspace{20mm}

A. The cabin measures more meters than before the ship accelerated.

B. The cabin measures fewer meters than before.

C. The cabin measures the same.

D. The effect is too small to be observable, because the ship's speed is still small compared
to the speed of light.
\end{frame}


\begin{frame}{Where u at?}
Use a spacetime diagram to resolve the following relativity paradox. Relativity
        says that in one frame of reference, event A could happen before event
        B, but in someone else's frame B would come before A. How can this be?
        Obviously the two people could meet up at A and talk as they cruised
        past each other. Wouldn't they have to agree on whether B had already
        happened?
\end{frame}

\begin{frame}{A relativistic cruise}

Liz has always been Grandpa Eddie's favorite grandchild, so he invites her
on an interstellar cruise, just like his own grandma did when he was little.
When Eddie was a kid, starships
only went
30\%
of the speed of light, but now they go 99.999999999\%
of $c$. As an experiment, Liz takes a flashlight and shines the beam along the ship's corridor
toward its front. Draw a spacetime diagram in the earth's frame of reference.
How does the beam's speed seem to Liz?

\end{frame}


\begin{frame}{Regions of spacetime}

\dq
The figure shows a spacetime diagram drawn in the standard way, with a single event A marked.

\vspace{3mm}

  \begin{mycolumns}

    \begin{column}{.3\textwidth}\includegraphics[width=3in]{ch03/figs/shade-spacetime}\end{column}

    \begin{column}{.7\textwidth}

    1. There is some set E of events that
    came \emph{earlier} in time than A. The set E fills some region of spacetime.
    With shading like this, \includegraphics{ch03/figs/shading-1}, indicate the region E.

   2. Similarly, use this kind of shading, \includegraphics{ch03/figs/shading-2}, to show the
   set C of events that could have \emph{caused} A.

    3. Which of these sets is a matter of opinion, depending on the frame of reference?

    \end{column}
  \end{mycolumns}


\end{frame}
